\nenumerisano{UVOD}
\label{ch:uvod}

\napomena{Kao u primjeru: UVOD je jedinstven tekst bez potpoglavlja, 4--5 strana.
Smjernice: definiše predmet i osnovne pojmove, daje teorijski okvir, ukazuje na nedostatke
postojećih rješenja, navodi istraživačka pitanja i predstavlja strukturu rada.
Teme redom (svaka po jedan ili više pasusa) --- tekst se može preuzeti iz prijave:}

\napomena{1. Kontekst i aktuelnost teme: mjesto asemblera u lancu kompajler → asembler →
linker → loader; širenje RISC-V arhitekture. \parencite{levine2000,hennessy2019}}

\napomena{2. Predmet istraživanja: dvoprolazni asembler za podskup RV32I.}

\napomena{3. Motivacija: razumijevanje veze hardvera i softvera, složenost industrijskih
alata (GNU as, LLVM), obrazovna vrijednost.}

\napomena{4. Cilj i svrha rada: opšti i posebni ciljevi (iz prijave).}

\napomena{5. Istraživačka pitanja: kao u primjeru --- „Prvo istraživačko pitanje glasi: …
Obrazloženje istraživačkog pitanja: …” za glavno i pomoćna pitanja.}

\napomena{6. Metode istraživanja: analiza i sinteza literature, modelovanje/specifikacija,
konstruktivna metoda, eksperimentalna (testiranje), komparativna, studija slučaja.}

\napomena{7. Struktura rada: po jedan pasus za poglavlja
\ref{ch:osnove}--\ref{ch:evaluacija}, zaključak i priloge.}
